% Generated draft: requires mathematical and source review.
\documentclass[11pt]{article}
\usepackage[margin=25mm]{geometry}
\usepackage{fontspec}
\setmainfont{Latin Modern Roman}
\usepackage{amsmath,amssymb,mathtools,unicode-math}
\setmathfont{Latin Modern Math}
\usepackage{enumitem,needspace,xurl,hyperref}
\setlength{\parindent}{0pt}
\setlength{\parskip}{5pt}
\setlength{\emergencystretch}{3em}
\newcommand{\N}{\mathbb{N}}
\newcommand{\Z}{\mathbb{Z}}
\newcommand{\Q}{\mathbb{Q}}
\newcommand{\R}{\mathbb{R}}
\newcommand{\C}{\mathbb{C}}
\newcommand{\F}{\mathbb{F}}
\newcommand{\A}{\mathbb{A}}
\newcommand{\PP}{\mathbb{P}}
\newcommand{\E}{\mathbb{E}}
\newcommand{\eps}{\varepsilon}
\newcommand{\dd}{\,\mathrm{d}}
\begin{document}
\section*{UnsolvedMath ID: 30001607}
\textbf{Source problem code:} catalogue entry
\subsection{Definitions and mathematical statement}
The supplied record combines an arrangement-complement conjecture with two random-complex threshold questions. Several conventions are missing; the formulations below make these qualifications explicit.

For a complex hyperplane arrangement \(\mathcal A\) in \(\C^\ell\), write
\[
M(\mathcal A)=\C^\ell\setminus\bigcup_{H\in\mathcal A}H.
\]
For a path-connected space \(X\), use the unreduced convention in which \(\operatorname{TC}(X)\) is the least number of open sets covering \(X\times X\) on each of which the endpoint map
\[
X^{[0,1]}\longrightarrow X\times X,\qquad \gamma\longmapsto(\gamma(0),\gamma(1)),
\]
admits a continuous section. This convention is suggested by the source's \(+1\), but is not explicitly specified there.

To explain the usual zero-divisor cup-length notation, let \(k\) be a coefficient field and let \(A=A^*(\mathcal A;k)\) denote the Orlik--Solomon algebra, identified in the usual finite-arrangement setting with \(H^*(M(\mathcal A);k)\). Give \(A\otimes_k A\) its graded tensor-product multiplication, and put
\[
\mu:A\otimes_k A\longrightarrow A,\qquad a\otimes b\longmapsto ab.
\]
The usual zero-divisor cup-length is
\[
\operatorname{zcl}_k(A)=\sup\left\{r\in\Z_{\geq0}:\ \exists z_1,\ldots,z_r\in\ker\mu,\ z_1\cdots z_r\ne0\right\}.
\]
The arrangement question asks for \(\operatorname{TC}(M(\mathcal A))\) and asks whether Yuzvinsky's conjectural equality, transcribed with the supplied notation,
\[
\operatorname{TC}(M(\mathcal A))
=\operatorname{zcl}_{A^*(\mathcal A)}(A^*(\mathcal A))+1,
\]
holds throughout the intended class of arrangements. If the displayed zero-divisor cup-length means the usual invariant just defined, its right-hand side is \(\operatorname{zcl}_k(A^*(\mathcal A;k))+1\). The record does not specify the coefficient field or ring, explain its subscript notation, or state whether only central arrangements are intended. These choices require confirmation before this becomes a fully specified universal conjecture.

For the random-complex questions, the usual two-dimensional Linial--Meshulam interpretation of \(Y(n,q)\) is the random simplicial complex on \(\{1,\ldots,n\}\), with every edge present and each triangle included independently with probability \(q\), and with no higher-dimensional simplices. The following formulation is conditional on this being the intended model. For \(n\geq3\) and fixed \(\alpha\leq0\), set \(q=n^\alpha\). The restriction \(\alpha\leq0\) is necessary for this expression to be a probability for arbitrarily large \(n\). Define the event
\[
E_H(n,\alpha)=\left\{H_1(Y(n,n^\alpha);\Z/p\Z)=0\ \text{for every prime }p\right\}.
\]
Here the quantifier over primes is simultaneous for each realized complex. Determine the threshold for \(E_H(n,\alpha)\) as \(n\to\infty\).

The second supplied event is the absence of an epimorphism from \(\pi_1(Y(n,n^\alpha))\) onto any finite group. Taken literally, this event is impossible, since every group surjects onto the trivial group. If the intended quantifier is over nontrivial finite groups, define
\[
E_F(n,\alpha)=\left\{\nexists\ \text{finite group }\Gamma\text{ with }|\Gamma|>1\ \text{and an epimorphism }\pi_1(Y(n,n^\alpha))\twoheadrightarrow\Gamma\right\}.
\]
The quantifier here is over all such finite groups simultaneously, rather than over one fixed group. Determine the threshold for this corrected event. The supplied record conjectures the threshold exponent \(\alpha=-\tfrac12\). If an exponent threshold is intended, this means
\[
\lim_{n\to\infty}\Pr(E_F(n,\alpha))=
\begin{cases}
0,&\alpha<-\tfrac12,\\
1,&-\tfrac12<\alpha\leq0,
\end{cases}
\]
for each fixed admissible \(\alpha\). The record does not specify behavior at \(\alpha=-\tfrac12\), a critical window, or whether a sharper notion of threshold is requested. The exclusion of the trivial group and this interpretation of threshold are explicit qualifications, not corrections verified against the report.
\subsection{Short English statement}
Determine the topological complexity of complex hyperplane arrangement complements and test the equality attributed to Yuzvinsky with zero-divisor cup-length plus one. Also determine when a Linial--Meshulam random complex simultaneously has vanishing first homology over every prime field, and when its fundamental group has no nontrivial finite quotient; the supplied conjectural exponent for the latter is \(-1/2\). Coefficients, arrangement scope, random-model conventions, and the exclusion of the trivial group require confirmation.
\subsection{Sources}
\begin{itemize}
\item Local UnsolvedMath registry: permanent ID \(30001607\), title ``Arrangement Complexity and Random-Complex Homology Thresholds,'' source problem code \texttt{catalogue entry}, collection ``Oberwolfach Reports Open Problems.'' The supplied record is the evidence used for this draft; local registry locations are \texttt{data/unsolvedmath/problems.json} and the supplemental \texttt{catalog/problems.json}.
\item Supplied report citation: \emph{Arbeitsgemeinschaft: Topological Robotics (2011)}, Oberwolfach Reports, pp.~2681--2744, DOI \texttt{10.4171/owr/2010/47}. Source URL: \url{https://doi.org/10.4171/owr/2010/47}. The record attributes the arrangement conjecture to Yuzvinsky and supplies no individual proposer for the combined entry. Its proposed year is \(2010\), while its report citation carries \(2011\); both are retained as supplied.
\item Supplied literature-source entry, labelled ``original Oberwolfach problem report'': \url{https://doi.org/10.4171/owr/2010/47}.
\item Literature assessment dated exactly \texttt{2026-08-21}: the supplied curation labels the entry open, reports results for special arrangements, and says that no general proof or counterexample for the combined claim was found. This is supplied dated curation, not an independently verified literature assessment. The background expressly requests independent verification of time-sensitive status and cited claims.
\item The registry labels the statement \texttt{corrected\_verified} and reports that its mathematical content was checked against the official report, with notation, context, and question form normalized. These are supplied provenance assertions; the report was not independently consulted for this draft.
\end{itemize}
\end{document}
